\section{Hesse型Poisson量子化の明示的な三次中心元}
\label{sec:explicit-hesse-center}
本節では量子化の定理が与える三次中心元を、滑らかなHesse型
三変数ブロックについて閉公式で指定する。
この公式はSklyanin代数の既知の中心三次式
\cite[Theorem 1およびSection 2]{DeLaetLeBruyn}
（原式は\cite[p.~35, Remark (1)]{SmithTate}）
のPoisson量子化パラメータへの代入である。
従って中心元の存在という既知の結果自体を新規な主張とはしない。
本研究で有用なのは、$0$を含む明示的なパラメータ開集合と
$\hbar=0$での留数を正規化した公式を同時に得ることである。

\subsection{公式と適用開集合}
$t\in\CC$、$t^3\ne-1$とし、
\begin{equation}
 f_t(x,y,z)=\frac{x^3+y^3+z^3}{3}+txyz
\end{equation}
と置く。$\hbar$を次数零のパラメータとして、
\begin{equation}
 [y,z]=\hbar\left(x^2+\frac{t}{2}(yz+zy)\right),
 \quad
 [z,x]=\hbar\left(y^2+\frac{t}{2}(zx+xz)\right),
 \quad
 [x,y]=\hbar\left(z^2+\frac{t}{2}(xy+yx)\right)
 \label{eq:hesse-commutators}
\end{equation}
による次数付き代数を$S_{\hbar,t}$と書く。
ここで
\begin{equation}
 S=x^3+y^3+z^3,\quad
 C=xyz+yzx+zxy,\quad
 D=yxz+xzy+zyx,\quad T=(C+D)/6.
\end{equation}
次の$\hbar$の多項式を定める：
\begin{equation}
 U(\hbar,t)=1+\frac{\hbar^2t^2}{2}
  +\frac{\hbar^4t}{6}-\frac{\hbar^4t^4}{48},
 \qquad
 W(\hbar,t)=
 \frac{(4-\hbar^2t^2)
 (12t+\hbar^2(t^3-8))}{48}.
 \label{eq:hesse-UW}
\end{equation}
また
\begin{equation}
 \Delta_{\hbar,t}=
 \left(t+\frac{\hbar^2}{3}
       \left(1+\frac{t^3}{4}\right)\right)^3
  +\left(1-\frac{\hbar^2t^2}{4}\right)^3.
 \label{eq:hesse-discriminant}
\end{equation}
と置く。$\Delta_{0,t}=1+t^3\ne0$である。

\begin{theorem}[Hesse中心三次式の閉公式]
\label{thm:explicit-hesse-center}
$t^3\ne-1$に対し、
\begin{equation}
 U_t=\{\hbar\in\mathbb A^1_{\CC}:
   U(\hbar,t)(1-\hbar^2t^2/4)\Delta_{\hbar,t}\ne0\}
\end{equation}
とする。$0\in U_t$であり、
各$\hbar\in U_t$に対して
$S_{\hbar,t}$は三生成元のKoszul AS-regular domainで、
Hilbert級数$(1-z)^{-3}$を持つ。
元
\begin{equation}
 \boxed{\quad
 \Omega_{\hbar,t}=
   \frac{U(\hbar,t)}3(x^3+y^3+z^3)
   +W(\hbar,t)\operatorname{sym}(xyz)
 \quad}
 \label{eq:hesse-central-element}
\end{equation}
は$S_{\hbar,t}$の次数三の中心正則元であり、
$\Omega_{0,t}=f_t$である。
特に$t=0$では
$\Delta_{\hbar,0}=1+\hbar^6/27$で、
\begin{equation}
 \Omega_{\hbar,0}
 =\frac{x^3+y^3+z^3}{3}
   -\frac{2\hbar^2}{3}\operatorname{sym}(xyz)
\end{equation}
を得る。
\end{theorem}
\begin{proof}
$q=\hbar t/2$、
$a=1-q$、$b=-(1+q)$、$c=-\hbar$と置く。
\eqref{eq:hesse-commutators}は
\begin{equation}
 ayz+bzy+cx^2=0,\quad
 azx+bxz+cy^2=0,\quad
 axy+byx+cz^2=0
\end{equation}
となる。$\hbar\ne0$かつ
$(1-\hbar^2t^2/4)\Delta_{\hbar,t}\ne0$なら、
$abc\ne0$で、標準Sklyaninの
\begin{equation}
 (a^3+b^3+c^3)^3\ne(3abc)^3
\end{equation}
が成り立つ。
実際、左辺の差は
$-27\hbar^3\Delta_{\hbar,t}$である。
従ってこの開部分では点曲線が滑らかな楕円曲線であり、
$S_{\hbar,t}$は三次元Sklyanin AS-regular domainである
\cite{SmithTate,DeLaetLeBruyn}。
$\hbar=0$の代数は可換多項式環である。

既知の中心元の公式は
\begin{equation}
 G=a(b^3-c^3)C+
   b(c^3-a^3)D+
   c(a^3-b^3)S
 \label{eq:hesse-sklyanin-G}
\end{equation}
である\cite[Section 2]{DeLaetLeBruyn}。
一方、三つの定義関係式にそれぞれ$x,y,z$を左から掛けて足すと
\begin{equation}
 C-D=\hbar S+\frac{\hbar t}{2}(C+D)
 \label{eq:hesse-CD}
\end{equation}
を得る。この等式で$G$を$S,C+D$で表し、
$a,b,c$を代入して係数を整理すると
\begin{equation}
 G=-9\hbar\left(
       \frac{U(\hbar,t)}3S+
       \frac{W(\hbar,t)}6(C+D)\right)
   =-9\hbar\Omega_{\hbar,t}.
 \label{eq:hesse-G-identity}
\end{equation}
これは$\hbar$で割る操作をせずに成立する多項式恒等式である。
$\hbar\ne0$の滑らかな楕円Sklyanin開部分では
$G$が中心なので$\Omega_{\hbar,t}$も中心である。
各次数成分が$\hbar$について有限自由であることは
三次元SklyaninのHilbert級数と零ファイバーにおける
同じ級数、および元論文のPBW構成から従う。
従って各交換子$[\Omega_{\hbar,t},x]$等は
有限自由加群の切断として稠密開集合上で零であり、
$U_t$の全体で零となる。
$\Omega_{0,t}=f_t$は$T_0=xyz$と
$U(0,t)=1$、$W(0,t)=t$から分かる。
既知のSklyanin中心元$G$は楕円開部分で非零であり、
代数はdomainだから正則である。
零ファイバーでも$f_t$は非零因子である。
\end{proof}


\begin{corollary}[任意の滑らかな三変数三次式]
\label{cor:arbitrary-smooth-ternary-cubic}
$g(x_1,x_2,x_3)$が滑らかな平面三次曲線を定める
斉次三次式であるとする。$g$のJacobian Poisson括弧
\begin{equation}
 [x_2,x_3]=\hbar\operatorname{sym}(\partial_{x_1}g),
 \quad [x_3,x_1]=\hbar\operatorname{sym}(\partial_{x_2}g),
 \quad [x_1,x_2]=\hbar\operatorname{sym}(\partial_{x_3}g)
\end{equation}
から得られる量子化は、$0$を含むZariski開集合上で
三次中心元の閉公式を持つ。
実際、Hesse標準形\cite[Lemma 1]{ArtebaniDolgachev}により
ある$M\in\mathrm{GL}_3(\CC)$、$\lambda\in\CC^\times$、
$t^3\ne-1$が存在して
\begin{equation}
 g(x)=\lambda f_t(Mx)
\end{equation}
となる。$y=Mx$、$\hbar'=\lambda\det(M)\hbar$とおけば、
\begin{equation}
 \boxed{\ \Omega^g_\hbar(x)=
  \lambda\Omega_{\hbar',t}(y)\big|_{y=Mx}\ }
 \label{eq:arbitrary-ternary-central}
\end{equation}
は三次中心元で、$\Omega^g_0=g$である。
\end{corollary}
\begin{proof}
Poissonの定数三重ベクトル
$\partial_{x_1}\wedge\partial_{x_2}\wedge\partial_{x_3}$
は座標変換$y=Mx$の下で
$\det(M)\,
 \partial_{y_1}\wedge\partial_{y_2}\wedge\partial_{y_3}$
に移る。従って$y$の交換子は
\begin{equation}
 [y_2,y_3]
 =\lambda\det(M)\hbar\,
   \operatorname{sym}(\partial_{y_1}f_t),
\end{equation}
およびその巡回置換となる。
対称化は線形座標変換と可換なので、量子化代数は
$S_{\hbar',t}$への次数付き同型を持つ。
その同型によって定理\ref{thm:explicit-hesse-center}の
中心元を戻し、留数が$g$になるよう$\lambda$倍する。
\end{proof}

\begin{corollary}[中心変数と独立ブロックの閉公式]
\label{cor:hesse-block-center}
相異なる三変数ブロック$x_s=(x_s,y_s,z_s)$と
中心変数$u$について
\begin{equation}
 F=\sum_{s=1}^k
 \left(\frac{x_s^3+y_s^3+z_s^3}{3}
       +t_sx_sy_sz_s\right)+P(u),
 \qquad t_s^3\ne-1,
\end{equation}
$P$は任意の斉次三次式とする。
各ブロックに\eqref{eq:hesse-commutators}を課し、
異なるブロック間と$u$を可換に取る。
各$U_{t_s}$の共通開部分で
\begin{equation}
 \Omega_\hbar
   =\sum_s\Omega_{\hbar,t_s}+P(u)
 \label{eq:hesse-multiblock}
\end{equation}
は次数三の中心元で、$\Omega_0=F$である。
各$u$の一次形式$\ell_{si}(u)$を選んで
$x_{si}$を$x_{si}+\ell_{si}(u)$に置き換えた
非自明な混合項を含むPoissonブロックでも、
同じ置換による明示公式が成立する。
\end{corollary}
\begin{proof}
各$S_{\hbar,t_s}$の中心元は他のブロックと可換であり、
$P(u)$も中心なので和は中心である。
中心線形シフトは生成元の次数付き同型を定め、
同型によって中心元が中心元へ移る。
\end{proof}


\subsection{可換復元の候補となる点曲線}
\begin{proposition}[楕円点曲線のHesseパラメータ]
\label{prop:hesse-point-cubic}
定理\ref{thm:explicit-hesse-center}の条件で$\hbar\ne0$とする。
$S_{\hbar,t}$のpoint schemeを定める滑らかな平面三次曲線は
\begin{equation}
 E_{\hbar,t}:\quad
 x^3+y^3+z^3+3\,\vartheta(\hbar,t)xyz=0,
\quad
 \boxed{\quad
 \vartheta(\hbar,t)=
 \frac{t+\frac{\hbar^2}{3}(1+t^3/4)}
 {1-\hbar^2t^2/4}
 \quad}
 \label{eq:hesse-point-parameter}
\end{equation}
である。$\vartheta(0,t)=t$と形式的に延長され、
\begin{equation}
 \Delta_{\hbar,t}
 =\left(1-\frac{\hbar^2t^2}{4}\right)^3
  \left(1+\vartheta(\hbar,t)^3\right).
 \label{eq:hesse-point-delta}
\end{equation}
特に$\hbar\in U_t$なら、右辺のHesse三次曲線は滑らかである。
\end{proposition}
\begin{proof}
滑らかな楕円型三次元Sklyanin代数のpoint schemeは
\begin{equation}
 abc(x^3+y^3+z^3)-(a^3+b^3+c^3)xyz=0
\end{equation}
である\cite{SmithTate,DeLaetLeBruyn}。
定理の証明で与えた
$a=1-\hbar t/2$、$b=-(1+\hbar t/2)$、$c=-\hbar$
について
\begin{equation}
 abc=\hbar(1-\hbar^2t^2/4),\qquad
 a^3+b^3+c^3
 =-3\hbar t-\hbar^3(1+t^3/4)
\end{equation}
を代入し、共通係数$\hbar$で割ると表示を得る。
最後の恒等式は\eqref{eq:hesse-discriminant}の定義から従う。
\end{proof}
\begin{remark}[幾何学的比較との関係]
この楕円点曲線と中心三次元から得られる非可換cubicの
rotation付き可換復元は論文Bで扱う。
具体的な一楕円ブロックの比較については
同論文の「一ブロックの明示的可換復元」を参照されたい。
ここでは中心三次量子化の存在・公式と圏同値を区別する。
\end{remark}

\begin{remark}[適用範囲と新規性]
\label{rem:hesse-scope}
一般の$f_s(x_s,u)$には$u$とactive変数の本質的な
混合項があり、必ずしも上の中心シフト形にはならない。
三変数制限が滑らかな場合、その全混合項を含む中心元は
系\ref{cor:pg-smooth-ternary-explicit}の有限多項式公式で得られる。
三変数制限が特異または零でも、全変数の三次超曲面が滑らかなら、
定理\ref{thm:pg-universal-determinant}が
原点で正則な有理式のCramer中心元を与える。
いずれの構成も形式的中心元の存在を入力としない。
また、\eqref{eq:hesse-central-element}は
周囲の量子化と中心元の明示式であり、
残余dg圏の相対saturationや
rotation付き可換復元を結論するものではない。
中心元の基本公式は古典的Sklyanin理論に由来し、
本節のHesseパラメータへの翻訳の新規性は未確認である。
\end{remark}
